% =====================================================================
\section{The Path Algebra of a Quiver}\label{sec:algebra}
% =====================================================================

This section is self-contained. We build, from first principles, the algebraic
objects that the model in Section~\ref{sec:model} relies on: quivers, paths, the
path algebra, its grading by length, the impact-degree neighbourhood system, and
the \emph{walk operators} that encode multi-hop reachability. We develop the
needed representation theory from first principles, assuming no prior exposure;
standard references are \cite{schiffler2014quiver,assem2006elements}.

% ---------------------------------------------------------------------
\subsection{Quivers and paths}
% ---------------------------------------------------------------------

\begin{definition}[Quiver]
A \emph{quiver} is a quadruple $Q=(Q_0,Q_1,s,t)$ where $Q_0$ is a finite set of
\emph{vertices}, $Q_1$ is a finite set of \emph{arrows}, and
$s,t\colon Q_1\to Q_0$ assign to each arrow $\alpha$ its \emph{source}
$s(\alpha)$ and \emph{target} $t(\alpha)$. We write $\alpha\colon i\to j$ when
$s(\alpha)=i$ and $t(\alpha)=j$.
\end{definition}

A quiver is thus a directed graph in which multiple arrows between the same
ordered pair of vertices are allowed (a \emph{multigraph}). For pattern
recognition this is exactly the data of a directed relational structure: vertices
are entities and arrows are directed relations, with parallel arrows recording
multiplicity. We write $n=|Q_0|$ and identify $Q_0$ with $\{1,\dots,n\}$.

\begin{definition}[Path]
A \emph{path of length} $k\ge 1$ in $Q$ is a sequence of arrows
$p=\alpha_k\cdots\alpha_2\alpha_1$ such that $t(\alpha_r)=s(\alpha_{r+1})$ for
$1\le r<k$; that is, the arrows compose head-to-tail. Its source is
$s(p)=s(\alpha_1)$ and its target is $t(p)=t(\alpha_k)$. For each vertex $i$ we
also admit a \emph{trivial path} $e_i$ of length $0$ with $s(e_i)=t(e_i)=i$.
\end{definition}

We read paths right-to-left, as for function composition. A path is a directed
\emph{walk}: vertices and arrows may repeat. The number of distinct walks of a
given length between two vertices is the quantity that drives over-squashing, and
the path algebra makes it algebraically explicit.

% ---------------------------------------------------------------------
\subsection{The path algebra and its grading}
% ---------------------------------------------------------------------

Let $\kk$ be a field (in practice $\kk=\RR$).

\begin{definition}[Path algebra]
The \emph{path algebra} $\kQ$ is the $\kk$-vector space with basis the set of all
paths of $Q$ (including the trivial paths $e_i$), endowed with the bilinear
multiplication given on basis paths by concatenation:
\[
  q\cdot p \;=\;
  \begin{cases}
    q\,p & \text{if } s(q)=t(p),\\[2pt]
    0 & \text{otherwise,}
  \end{cases}
\]
and extended linearly. The trivial paths satisfy $e_{t(p)}\,p=p=p\,e_{s(p)}$ and
$e_ie_j=\delta_{ij}e_i$, so $\{e_i\}_{i\in Q_0}$ is a complete set of orthogonal
idempotents and $1=\sum_{i}e_i$ is the unit.
\end{definition}

The multiplication respects path length: concatenating a path of length $k$ with
one of length $\ell$ yields a path of length $k+\ell$ (or zero). Hence $\kQ$ is a
\emph{graded} algebra,
\[
  \kQ \;=\; \bigoplus_{k\ge 0}\kQ_k,
  \qquad \kQ_k\cdot \kQ_\ell \subseteq \kQ_{k+\ell},
\]
where $\kQ_k$ is the subspace spanned by the paths of length exactly $k$. In
particular $\kQ_0=\bigoplus_i \kk e_i$ and $\kQ_1$ is spanned by the arrows.

The idempotents $e_i$ let us isolate the paths between a fixed pair of vertices.
For $i,j\in Q_0$ and $k\ge 0$, the subspace
\[
  e_j\cdot \kQ_k\cdot e_i \;\subseteq\; \kQ
\]
is spanned exactly by the paths of length $k$ from $i$ to $j$. Its dimension is
therefore the \emph{number of directed walks} of length $k$ from $i$ to $j$. The
following elementary proposition states the correspondence with linear algebra
precisely; it underlies the model's computation.

\begin{proposition}[Graded multiplicity equals walk count]\label{prop:count}
Let $\Adj\in\NN^{n\times n}$ be the adjacency matrix of $Q$, $\Adj_{ij}=\#\{\alpha
\in Q_1: \alpha\colon i\to j\}$. Then for all $k\ge 0$ and all $i,j\in Q_0$,
\[
  \dim_{\kk}\bigl(e_j\cdot \kQ_k\cdot e_i\bigr)
  \;=\; (\Adj^{\,k})_{ij}.
\]
\end{proposition}

\begin{proof}
By induction on $k$. For $k=0$, $e_j\kQ_0 e_i=\delta_{ij}\,\kk e_i$, of dimension
$\delta_{ij}=(\Adj^0)_{ij}=(\mathrm{Id})_{ij}$. For $k=1$, a basis of
$e_j\kQ_1 e_i$ is the set of arrows $i\to j$, of cardinality $\Adj_{ij}$.
Assume the claim for $k-1$. Every path of length $k$ from $i$ to $j$ factors
uniquely as $\alpha\cdot p$, where $p$ is a path of length $k-1$ from $i$ to some
intermediate vertex $r$ and $\alpha\colon r\to j$ is an arrow. The number of such
factorisations is $\sum_{r}\Adj_{rj}\,(\Adj^{\,k-1})_{ir}=(\Adj^{\,k})_{ij}$,
which proves the claim. \qed
\end{proof}

We call the integer $n_k(i,j):=(\Adj^{\,k})_{ij}=\dim_\kk(e_j\kQ_k e_i)$ the
\emph{walk multiplicity} from $i$ to $j$ at range $k$. It is exactly the number
of length-$k$ messages that a message-passing GNN forces through node $j$ from
node $i$ after $k$ layers, and hence the source of over-squashing.

% ---------------------------------------------------------------------
\subsection{Impact degree and the fundamental neighbourhood system}
\label{sec:sfv}
% ---------------------------------------------------------------------

The grading by length induces a canonical stratification of the vertices around
each node, which organises the multi-hop structure the model will exploit. For a
quiver $Q$, the \emph{impact degree} $g(i,j)$ is the length of the shortest
directed path from $i$ to $j$ (with $g(i,i)=0$ and $g(i,j)=\infty$ if no directed
path exists); equivalently, $g(i,j)=\min\{k: (\Adj^{\,k})_{ij}>0\}$, the first
range at which the walk multiplicity becomes positive.

\begin{definition}[Fundamental neighbourhood system]\label{def:sfv}
For a node $c\in Q_0$, its \emph{in-neighbourhood} of radius $k$ is
\[
  A_k(c) \;=\; \{\,i\in Q_0 : g(i,c)\le k\,\},
\]
the set of nodes that can send influence to $c$ along a directed path of length
at most $k$. The nested family $\{A_k(c)\}_{k\ge 0}$, with
$A_0(c)=\{c\}\subseteq A_1(c)\subseteq\cdots$, is the \emph{fundamental
neighbourhood system} (FNS) of $c$, and the layer
$\mathcal{N}_k(c)=A_k(c)\setminus A_{k-1}(c)$ collects the nodes at impact degree
exactly $k$.
\end{definition}

The FNS is the combinatorial object that the walk operators of
Definition~\ref{def:walkop} act on. The precise relation between its layers and
walk-reachability is the following.

\begin{proposition}[FNS layers vs.\ walk reachability]\label{prop:fns}
For $k\ge1$ let $\mathcal{S}_k=\{(c,i):(\Adj^{\,k})_{ic}>0\}$ be the pairs
joined by a length-$k$ walk $i\walk c$. Then
\[
  \{(c,i): i\in \mathcal{N}_k(c)\} \;\subseteq\; \mathcal{S}_k ,
\]
and the inclusion is strict in general: a length-$k$ walk may also reach nodes
at impact degree $<k$. If $Q$ is \emph{graded}---there is
$\rho\colon Q_0\to\mathbb{Z}$ with $\rho(t(\alpha))=\rho(s(\alpha))+1$ for every
arrow $\alpha$---then equality holds for every $k$.
\end{proposition}

\begin{proof}
If $i\in\mathcal{N}_k(c)$ then $g(i,c)=k$ and a shortest path is a length-$k$
walk, so $(\Adj^{\,k})_{ic}>0$. Strictness: in the quiver $1\to2\to3$ with an
extra arrow $1\to3$ one has $(\Adj^{2})_{13}>0$ yet $g(1,3)=1$. If $Q$ is
graded, every walk $i\walk c$ has length $\rho(c)-\rho(i)$; a length-$k$ walk
forces $\rho(c)-\rho(i)=k$, so every walk---in particular a shortest
path---has length $k$, giving $g(i,c)=k$ and $i\in\mathcal{N}_k(c)$. \qed
\end{proof}

The benchmark graphs of Section~\ref{sec:exp} are graded (the layer index is
such a $\rho$), so there the supports coincide with the FNS layers; on graphs
with cycles, such as the ring, $\mathcal{S}_k$ strictly contains the layer
relation. Stacking ranges $k=1,2,\dots$ therefore walks the embedding outward
through the layers of the FNS, revisiting earlier layers where the topology
permits. The multiplicity $n_k(i,c)$ records \emph{how many} length-$k$ paths
populate each layer---the textured, algebraic refinement of the scalar impact
degree.

% ---------------------------------------------------------------------
\subsection{An explicit example}
% ---------------------------------------------------------------------

\begin{example}[A diamond quiver]\label{ex:diamond}
Let $Q$ have vertices $\{1,2,3,4\}$ and arrows
$\alpha_1\colon 1\to2$, $\alpha_2\colon 1\to3$, $\beta_1\colon 2\to4$,
$\beta_2\colon 3\to4$ (Figure~\ref{fig:diamond}). Its adjacency matrix and the
square are
\[
  \Adj=\begin{pmatrix}0&1&1&0\\0&0&0&1\\0&0&0&1\\0&0&0&0\end{pmatrix},
  \qquad
  \Adj^2=\begin{pmatrix}0&0&0&2\\0&0&0&0\\0&0&0&0\\0&0&0&0\end{pmatrix}.
\]
There are two length-$2$ walks from $1$ to $4$, namely
$\beta_1\alpha_1$ and $\beta_2\alpha_2$, so
$\dim(e_4\kQ_2 e_1)=(\Adj^2)_{14}=2$, in agreement with
Proposition~\ref{prop:count}. These two walks are \emph{parallel}: same source,
same target, same length. Whether they should be treated as one channel or two is
a modelling decision we return to in Section~\ref{sec:exp}.
\end{example}

\begin{figure}[t]
\centering
\begin{tikzpicture}[->,>={Stealth[round]},node distance=16mm and 18mm,
  every node/.style={circle,draw,minimum size=7mm,inner sep=0pt,font=\small}]
  \node (1) {1};
  \node (2) [above right=of 1] {2};
  \node (3) [below right=of 1] {3};
  \node (4) [below right=of 2] {4};
  \draw (1) -- node[draw=none,fill=none,above left,font=\scriptsize]{$\alpha_1$} (2);
  \draw (1) -- node[draw=none,fill=none,below left,font=\scriptsize]{$\alpha_2$} (3);
  \draw (2) -- node[draw=none,fill=none,above right,font=\scriptsize]{$\beta_1$} (4);
  \draw (3) -- node[draw=none,fill=none,below right,font=\scriptsize]{$\beta_2$} (4);
\end{tikzpicture}
\caption{The diamond quiver of Example~\ref{ex:diamond}. The two length-$2$ walks
$1\walk 4$ are parallel; $\dim(e_4\kQ_2 e_1)=2$.}
\label{fig:diamond}
\end{figure}

% ---------------------------------------------------------------------
\subsection{Walk operators}
% ---------------------------------------------------------------------

Proposition~\ref{prop:count} lets us turn the graded path algebra into a family of
linear operators on node features, one per range $k$. These \emph{walk operators}
are the bridge to the neural model.

\begin{definition}[Walk operator]\label{def:walkop}
For each range $k\ge 1$ define the matrix $W^{(k)}\in\RR^{n\times n}$ by
\[
  W^{(k)}_{j i}\;=\;\frac{(\Adj^{\,k})_{ij}}{\,Z_k\,},
\]
where $Z_k=\max_{j}\sum_i (\Adj^{\,k})_{ij}$ is a single normalising constant for
the range. The \emph{walk-reachability support} at range $k$ is the index set
\[
  \mathcal{S}_k \;=\; \bigl\{(j,i): (\Adj^{\,k})_{ij}>0\bigr\}
  \;=\;\bigl\{(j,i): \dim(e_j\kQ_k e_i)>0\bigr\}.
\]
\end{definition}

Acting on a feature matrix $H\in\RR^{n\times d}$, the operator $W^{(k)}H$ sends to
each node $j$ a weighted sum of the features of every node $i$ that reaches $j$ by
a length-$k$ walk, the weight being proportional to the walk multiplicity
$n_k(i,j)$. Two facts are worth stressing.

\emph{(i) Global normalisation, not row normalisation.} We divide by a single
constant $Z_k$ per range rather than normalising each row to sum to one. Row
normalisation would map the multiplicities to a uniform distribution over the
reachable sources and would erase the very multiplicity information that
distinguishes $W^{(k)}$ from a plain reachability mask; the global constant keeps
the relative multiplicities intact while bounding the magnitude.

\emph{(ii) Sparsity.} The support $\mathcal{S}_k$ is exactly the nonzero pattern
of $\Adj^{\,k}$. On the structured graphs of interest it is a small fraction of
all $n^2$ pairs (see Section~\ref{sec:exp}), so $W^{(k)}$ can be stored and
applied in $O(|\mathcal{S}_k|)$ rather than $O(n^2)$.

The fixed-weight operator $W^{(k)}$ already mitigates over-squashing, because it
lets node $j$ read from range-$k$ sources \emph{directly}, in one step, instead
of forcing the signal through $k$ rounds of local message passing. Its
limitation---which motivates Walk Attention---is that the weights are fixed by
multiplicity and cannot \emph{select} which reachable source matters. We address
this in Section~\ref{sec:model} by learning the weights over the same support.

% ---------------------------------------------------------------------
\subsection{The quotient algebra}
% ---------------------------------------------------------------------

Representation theory routinely passes from $\kQ$ to a quotient $\kQ/I$ by an
\emph{admissible ideal} $I$, identifying paths that are declared equal. In the
diamond of Example~\ref{ex:diamond}, the relation $\beta_1\alpha_1\equiv
\beta_2\alpha_2$ generates an ideal $I$ for which
$\dim\bigl(e_4(\kQ/I)_2 e_1\bigr)=1$: the two parallel walks collapse to a single
class. More generally, $\dim\bigl(e_j(\kQ/I)_k e_i\bigr)\le n_k(i,j)$, so the
quotient yields a structurally smaller multiplicity.

It is tempting to use this quotient directly as an over-squashing remedy:
replace the walk multiplicities $n_k(i,j)$ by the smaller quotient dimensions,
thereby ``de-duplicating'' redundant paths before aggregation. We tested this and
found it \emph{counterproductive} (Section~\ref{sec:exp}): collapsing parallel
walks destroys multiplicity information that downstream tasks may require. The
constructive role of the quotient, then, is not to compress the signal but to
\emph{reshape the support}---e.g.\ to route distinct path classes to distinct
attention heads---a direction we leave to future work. In the model that follows
we use the path algebra only through the reachability support $\mathcal{S}_k$ and
the (unquotiented) walk operators $W^{(k)}$.

% ---------------------------------------------------------------------
\subsection{A common algebraic ground}\label{sec:ground}
% ---------------------------------------------------------------------

The objects above---the graded path algebra $\kQ$, the impact degree and its
fundamental neighbourhood system (FNS), and the quotient $\kQ/I$---form one
algebraic substrate that specialises in two directions, according to how a
node-update map $\varphi$ over the layers $\mathcal{N}_k(c)$ is read. In the
\emph{dynamical} reading, $\varphi$ is a fixed evolution rule applied over the
FNS, weighting influence by impact degree; iterating it in time yields the
\emph{automata of impact over quivers} (AIQ) framework, with its SI/SIS/SIR rules
and category structure, developed separately
\cite{zainea2026aiq,ibanez2025casimulations}. In the \emph{representational}
reading, taken here, $\varphi$ is \emph{learned}: a vector space at each vertex
and a linear map at each arrow is a representation of $Q$ (a $\kQ$-module, via
$\mathrm{Rep}(Q)\cong\kQ\text{-}\mathrm{Mod}$
\cite{schiffler2014quiver,assem2006elements}), and a graph neural network on $Q$
is such a representation with learned arrow maps. Walk Attention is its graded
instance: a learned representation whose support is the FNS. The two readings
share the trunk and differ only in whether $\varphi$ is prescribed or learned.
